% ===================== 第 7 章 =====================
\section{证据门、三层审计与证据谱系}

\subsection{六阶段证据门}

每个候选事件在 \code{packages/evidence/gate.py} 的 \code{FullEvidenceGate} 中依次通过六个阶段，前一阶段不通过就不进入下一阶段，审计结论逐行落库。

\begin{table}[ht]
\centering\small
\caption{证据门六个阶段的职责}
\label{tab:gate}
\begin{tabularx}{\linewidth}{@{}L{4.0cm} X@{}}
\toprule
\rowcolor{accent}
\tblhead 阶段 & \tblhead 检查内容 \\
\midrule
SCHEMA              & 字段完整性与类型正确性 \\
DEDUPLICATION       & 与既有对象的语义重复判定，重复即不消耗后续预算 \\
SOURCE              & DOI 解析、标题/作者/年份一致性、撤稿状态 \\
CITATION\_ENTAILMENT& 原文是否真的支持被绑定的主张 \\
METHOD\_QUALITY     & 直接性、设计强度、测量效度、统计精度、可复现性、外部效度 \\
GRAPH\_CONSISTENCY  & 是否制造矛盾节点、重复分叉谱系等图约束破坏 \\
\bottomrule
\end{tabularx}
\end{table}

第六阶段不是恒真空检查：它调用 \code{consistency.py::check\_graph\_consistency()}，所需布尔值来自 \code{SqlGraphConsistencyQuery} 对数据库的真实查询。第三、五阶段读取候选事件自带的元数据，撤稿论文在此被拒。

\subsection{证据分级 A–D：等级不够就进不了结论}

\begin{table}[ht]
\centering\small
\begin{tabularx}{\linewidth}{@{}c L{4.6cm} X@{}}
\toprule
\rowcolor{accent}
\tblhead 等级 & \tblhead 可得性 & \tblhead 处置 \\
\midrule
A & 全文与精确原文可得      & \code{ADMIT}，节点状态 active \\
B & 仅有摘要与可靠元数据    & \code{SOURCE\_ONLY}，节点 provisional，不得单独支撑高置信因果结论 \\
C & 综述中的二手描述        & \code{DISCOVERY\_ONLY}，留账本、不入图 \\
D & 网页或新闻线索          & \code{TOOL\_LEAD\_ONLY}，留账本、不入图 \\
\bottomrule
\end{tabularx}
\end{table}

\begin{honestbox}
\boxtag{\color{warnamber}实事求是：当前管线只产 B 级}
首版取证管线能稳定拿到的是题录与摘要元数据，因此正式发现最高只到 B 级。这不是配置遗漏：在拿不到全文时把证据标成 A 级，等于默许高置信因果结论建立在摘要之上。等级标签是承诺，不是装饰。
\end{honestbox}

\subsection{三层审计：防止六类常见误写}

正式 Finding 必须依次通过来源真实性（\code{source\_verifier.py}）、引用蕴含（\code{citation\_verifier.py}）与方法质量（\code{method\_auditor.py}）三层审计。引用蕴含层专门拦截六类高频误写：引言里的假设被写成实测结果；讨论段的推测被写成实证发现；不显著被写成「没有效应」；相关被升级为因果；子群结果被外推为总体结论；作者未检验的机制被写成已验证机制。

\subsection{因果越级拦截}

\code{CausalUpgradePolicy} 在证据门内专门识别「横截面设计 + 因果主张」的组合，命中即隔离事件并留审计记录。「相关不等于因果」在多数系统里是一句提示语，在这里是一条不可绕过的判定规则。

\subsection{证据谱系：论文数量不等于独立证据数量}

\code{lineage.py} 定义八类证据依赖关系：同一数据集、重叠样本、同一研究团队、预印本与正式版、扩展研究、再分析、引用但无新数据、综述纳入。\code{independence.py} 的 \code{cluster\_evidence} 用并查集把存在依赖的研究聚成独立证据簇——同一实验室的两个不同数据集不合并，但共享同一份数据的六篇论文只算一簇。界面与报告会同时展示论文数与独立证据簇数，两个数字的差距本身就是在向研究者报警。

\subsection{防止共享错误被当成共识的四道防线}

\begin{enumerate}[leftmargin=1.6em, itemsep=3pt]
\item \textbf{盲证据评审。}首轮判断剥离作者、期刊、被引量等声誉信号，只给方法、样本与结果。当前由结构性测试锁定网关不透传声誉字段；主动剥离的轮次处理器属保留项（第 12 章）。
\item \textbf{独立双抽取。}高影响论文由两个不同模型/角色独立抽取后比对。首版未实现，列入路线图。
\item \textbf{来源多样性约束。}当一项主张的全部支持证据来自同一数据集、团队、国家、测量方式或设计时，\code{source\_diversity.py} 自动生成盲点，并禁止给出高置信结论。
\item \textbf{对抗式检索。}\code{adversarial\_retrieval.py} 为每条已确认主张生成六类反向检索意图（反驳、零结果、替代理论、测量批评、复现失败、边界翻转）并发起真实查询；查不到就如实记录尝试次数，不伪造命中。
\end{enumerate}
